% Generated by scripts/assemble.py; corrections belong in source/corrections.json.
\examhead{137}{武汉理工大学 2026 年数学分析试卷}{130}
% p131-b23; source PDF page 131
\par\Needspace{5\baselineskip}
\blockanchor{p131-b23}
\majorquestion{一}计算题. 每题 10 分, 共 60 分.\par
% p131-b24; source PDF page 131
\blockanchor{p131-b24}
\question{1}求极限 $\lim_{n\to \infty}n^2\left(\arctan{\frac{2026}{n}} -\arctan{\frac{2026}{n + 1}}\right)$\par
% p132-b00; source PDF page 132
\blockanchor{p132-b00}
\question{2}设 $z = z(x, y)$ 由 $f(z - x, z - y) = 0$ 所确定, 其中 $f(u, v)$ 具有连续的一阶偏导数, 且 $\frac{\partial f}{\partial u} + \frac{\partial f}{\partial v} \neq 0$, 求 $\frac{\partial z}{\partial x} + \frac{\partial z}{\partial y}$.\par
% p132-b01; source PDF page 132
\blockanchor{p132-b01}
\question{3}已知 $f(x) = \int_{\sqrt{x}}^{\sqrt{\frac{\pi}{2}}} \frac{1}{1 + \tan y^2} \, \mathrm{d}y$，求 $\int_0^{\frac{\pi}{2}} \frac{f(x)}{\sqrt{x}} \, \mathrm{d}x$.\par
% p132-b02; source PDF page 132
\blockanchor{p132-b02}
\question{4}求第二型曲面积分 $\iint_{S} \frac{x \, \mathrm{d}y \, \mathrm{d}z + z \, \mathrm{d}x \, \mathrm{d}y}{\sqrt{(x^2 + y^2 + z^2)^3}}$, 其中 $S$ 为 $z = -\sqrt{a^2 - x^2 - y^2}$, 方向向上.\par
% p132-b03; source PDF page 132
\blockanchor{p132-b03}
\question{5}求由方程 $x^{2} + 2xy + 2y^{2} = 1$ 确定的隐函数 $y = y(x)$ 的所有极值, 并判断是极大值还是极小值.\par
% p132-b04; source PDF page 132
\blockanchor{p132-b04}
\question{6}求幂级数 $\sum_{n=1}^{\infty} (-1)^{n-1} n(x - 1)^n$ 的收敛域与和函数.\par
% p132-b05; source PDF page 132
\par\Needspace{5\baselineskip}
\blockanchor{p132-b05}
\majorquestion{二}解答题. 每题 15 分, 共 30 分.\par
% p132-b06; source PDF page 132
\par\Needspace{6\baselineskip}
\blockanchor{p132-b06}
\question{1}已知二元函数\par
\nopagebreak[4]
% p132-b07; source PDF page 132
\blockanchor{p132-b07}
\displayformula{f (x, y) = \left\{ \begin{array}{l l} \frac {x ^ {3} + x ^ {2} y + x y ^ {2} + 2 y ^ {4}}{x ^ {2} + y ^ {2}}, & x ^ {2} + y ^ {2} \neq 0; \\ 0, & x ^ {2} + y ^ {2} = 0. \end{array} \right.}
% p132-b08; source PDF page 132
\blockanchor{p132-b08}
(1) 判断 $f(x, y)$ 在 $(0, 0)$ 处的连续性.\par
% p132-b09; source PDF page 132
\blockanchor{p132-b09}
(2) 判断  $f(x,y)$  在  $(0,0)$  处的两个一阶偏导数是否存在？若存在请求其值.\par
% p132-b10; source PDF page 132
\blockanchor{p132-b10}
(3) 判断 $f(x, y)$ 在 $(0, 0)$ 处的可微性, 请详细说明理由.\par
% p132-b11; source PDF page 132
\blockanchor{p132-b11}
\question{2}当 $p > 0$ 时, 讨论级数 $\sum_{n=2}^{\infty} \ln \left(1 + \frac{(-1)^n}{n^p}\right)$ 在 $p$ 为何值时绝对收敛, 条件收敛及发散.\par
% p132-b12; source PDF page 132
\par\Needspace{5\baselineskip}
\blockanchor{p132-b12}
\majorquestion{三}证明题. 每题 15 分, 共 60 分.\par
% p132-b13; source PDF page 132
\blockanchor{p132-b13}
\question{1}证明：函数列 $f_{n}(x) = \frac{\ln(nx)}{nx^{2}}$ 在 $(1, +\infty)$ 上一致收敛.\par
% p132-b14; source PDF page 132
\blockanchor{p132-b14}
\question{2}证明: $f(x)$ 在 $(0,1)$ 上一致连续的充分必要条件为对任意 $(0,1)$ 中的柯西列 $\{x_{n}\}, \{f(x_{n})\}$ 也为柯西列.\par
% p132-b15; source PDF page 132
\blockanchor{p132-b15}
\question{3}证明: 数列 $a_{n} = 1 + \frac{1}{\sqrt{2}} + \frac{1}{\sqrt{3}} + \cdots + \frac{1}{\sqrt{n}} - 2\sqrt{n}$ 收敛.\par
% p132-b16; source PDF page 132
\blockanchor{p132-b16}
\question{4}解答如下问题:\par
% p132-b17; source PDF page 132
\blockanchor{p132-b17}
(1) 设 $k$ 为正整数, 证明: $\int_{(k-1)\pi}^{k\pi} \frac{1}{1 + x^2 \sin^2 x} \mathrm{d}x \geq 2 \int_0^{\frac{\pi}{2}} \frac{1}{1 + k^2 \pi^2 \sin^2 x} \mathrm{d}x.$\par
% p132-b18; source PDF page 132
\blockanchor{p132-b18}
(2) 证明: 反常积分 $\int_{0}^{+\infty} \frac{1}{1 + x^2 \sin^2 x} \, \mathrm{d}x$ 发散.\par
